% !PRINTED 113
\weberpage{113}
%<<CONT
und die durch die Vertauschung von $i$ mit $-i$ daraus hervorgehende
 Zahl, ferner mit $c$ eine rationale Zahl bezeichnen, setzen
\[
\begin{aligned}
a + b i &= i^{\lambda}\, c^2 H_1^4\\
a - b i &= i^{-\lambda} c^2 H_2^4.
\end{aligned}
\]
Aus (8) folgt durch Ausziehen der $4^{\text{ten}}$ Wurzel, wodurch
eine Potenz von $i$ als Factor hinzukommen kann,
\weq{15}{\begin{aligned}
(i, x_0) &= i^h \sqrt[4]{i}^{\lambda} \sqrt{c}\, H_1\\
(-i, x_0) &= i^{-h} \sqrt[4]{-i}^{\lambda} \sqrt{c}\, H_2\\
(-1, x_0) &= \sqrt{m},
\end{aligned}}
und daraus ergiebt sich, da $\sqrt[4]{i}$, $\sqrt{c}$ und $\sqrt{m}$ Kreistheilungszahlen
 sind, die Richtigkeit unseres Satzes auch in diesem Falle.

Auch hier kann die weitere Betrachtung des Baues der in
(15) vorkommenden Ausdr\"ucke dazu dienen, die Zusammensetzung
 der $x_0, x_1, x_2, x_3$ durch die Perioden $\eta$ n\"aher zu erforschen,
was aber wieder zu keinen anderen Resultaten f\"uhren
kann, als zu den schon in \S.\,24 abgeleiteten.

Wir haben also hiermit den Satz allgemein bewiesen, \gesperrt{dass
alle Abel'schen K\"orper dritten und vierten Grades
Kreistheilungsk\"orper sind, und dass alle diese K\"orper
rational durch die Kreistheilungsperioden dargestellt
werden k\"onnen.}

\rule{3cm}{0.4pt}
